% ============================================================
%  CHAPTER 1 — Project Foundations: Context, Analysis, and Methodology
% ============================================================
\chapter{Project Foundations: Context, Analysis, and Methodology}

% ── 1.1 Introduction ─────────────────────────────────────────
\section{Introduction}

This chapter presents the foundations of the Sahhty project by describing the environment
in which it was carried out, the problem it was designed to solve, and the approach that
guided its development from beginning to end.

The chapter opens with a presentation of the organization that hosted this project, followed
by an overview of the subject and the motivations behind it. A problem statement identifies
the specific difficulties faced by patients, and particularly by pregnant women, in managing
their health in the current healthcare landscape. A review of existing solutions examines
what the market currently offers and highlights the limitations that justify building something
new. The proposed solution is then introduced, along with the development approach and
project timeline adopted by the team.

% ── 1.2 Presentation of the Host Organization ────────────────
\section{Presentation of the Host Organization}

This project was carried out as a final-year internship hosted at \textbf{Pharmacie Nahla
Noureddine}, a community pharmacy located in Tunisia. The project originated from an
initiative taken by the pharmacy itself: observing a real and recurring need among its
patients, the pharmacy approached the development team and proposed the creation of a
digital health companion tailored to the specific challenges faced by pregnant women and
patients in the local healthcare context. The team accepted this proposal and worked
closely with the pharmacy throughout the entire project, drawing on its professional
expertise to shape the priorities and clinical direction of the application.

Pharmacies occupy a central position in the daily healthcare journey of patients. Beyond
dispensing medications, they serve as a first point of contact for health advice, drug
interaction warnings, and guidance on medication safety during pregnancy. This direct
contact with patients and their everyday health concerns gave the pharmacy a clear and
concrete understanding of the problem to be solved, making it both the initiator and the
primary source of domain knowledge for this project.



% ── 1.3 Subject Presentation ─────────────────────────────────
\section{Subject Presentation}

The subject of this project is the development of a mobile health platform
called \textbf{Sahhty}, a word meaning ``my health'' in Tunisian Arabic. The platform
serves two types of users: patients (both pregnant women and men who want to monitor
their general health) and doctors who wish to manage their patients and appointments
through a dedicated interface.

The application was born from a straightforward observation. Despite the rapid growth of
health applications worldwide, pregnant women in Tunisia still lack a single, unified tool
that covers all aspects of their pregnancy journey. They record their measurements in
paper notebooks, carry heavy medical folders from one appointment to the next, and have
no reliable way to check whether a medication is safe to take during their current stage of
pregnancy.

Sahhty was designed to address all of these needs in one place. It allows patients to record
and track their health readings over time, follow the progress of their pregnancy week by
week, search for doctors near them and book appointments, store and share their medical
documents securely, look up the safety of medications during pregnancy, and receive alerts
when something requires their attention.

For doctors, the platform provides a space to manage their schedule, review the files of
patients who have granted them access, respond to appointment requests, and stay informed
about their practice activity through real-time notifications.

The platform is designed exclusively for Android smartphones and is available in both
French and English, with the interface adapting its colours and content depending on
whether the user is a woman or a man.

% ── 1.4 Problem Statement ─────────────────────────────────────
\section{Problem Statement}

\subsection*{The Growth of Health Applications}

Since the Internet spread to virtually every aspect of daily life, health-related applications
have grown at an extraordinary pace. According to the World Health Organization, the use
of mobile technologies for health purposes gives users access to information, services, and
guidance, and contributes to meaningful improvements in health behaviours that help
prevent both acute and chronic conditions~\cite{who2016}. The scale of this shift is
considerable: the WHO reports more than seven billion mobile phone subscriptions
worldwide, and in some countries, access to a mobile phone is easier than access to clean
water or electricity~\cite{who2018}.

Health applications installed on smartphones and wearable devices have emerged as one of
the most direct ways for people to take an active role in managing their own health. Some
collect information automatically without any input from the user, monitoring steps, heart
rate, or sleep patterns throughout the day. Others ask the user to enter information
themselves, such as their diet, medication schedule, or clinical measurements. The
combination of both approaches makes it possible for these tools to suggest changes in
daily habits and reduce the risk of avoidable health problems~\cite{milne2020}.

\subsection*{The Fragmented Reality of Healthcare in Tunisia}

Alongside consumer health applications, a second category of digital health tools exists
within hospitals and healthcare institutions. These systems fully digitise the patient record
and allow every practitioner within the institution to access the patient's complete history,
test results, prescriptions, and ongoing treatments. In such environments, doctors can avoid
repeating examinations that have already been performed and can detect dangerous
combinations of medications before they are prescribed, because all the information is
visible and centralised.

In the Tunisian healthcare system, this level of organisation is largely absent. Patient
records are not systematically computerised, and there is little or no structured
communication between different healthcare facilities. As a result, patients often find
themselves carrying heavy paper folders containing a disorganised mix of laboratory results,
radiological reports, and prescriptions issued by different doctors at different locations.

This fragmentation creates serious consequences. At best, it leads to unnecessary
repetition: the same examination is ordered multiple times because no one can see that it
was already performed. At worst, it creates dangerous situations where two medications
prescribed by two different doctors interact with each other in a harmful way, and no one
is aware of it because neither doctor has access to the other's prescription.

\subsection*{A More Serious Problem for Pregnant Women}

While this situation affects all patients, it becomes particularly concerning for pregnant
women, because every health decision made during pregnancy affects not only the mother
but also the child she is carrying. A medication that is safe for a healthy adult may cause
serious harm to a developing baby at a specific stage of pregnancy. A drug combination
that produces only mild discomfort in a non-pregnant person can have lasting consequences
for fetal development.

pregnant women in Tunisia therefore face a dual challenge. They must navigate a
fragmented healthcare system while also managing a level of medical complexity that
requires precise, personalised information about the safety of every treatment they take.
Fragmented records, multiple prescribers, and the absence of any reliable way to check
medication safety during pregnancy together define the core problem that this project
was created to address.


% ── 1.5 Study of Existing Solutions ──────────────────────────
\section{Study of Existing Solutions}

\subsection{Market Overview}

The global market for health applications has expanded significantly over the past decade.
Four main categories of solutions are currently available to patients:

\textbf{Cycle and fertility tracking applications} help women monitor their menstrual cycle,
predict ovulation, and track the early stages of pregnancy. They are widely used and
scientifically grounded, but they focus exclusively on the reproductive dimension of health
and do not address broader health monitoring or clinical risk evaluation.

\textbf{Pregnancy companion applications} provide week-by-week information about fetal
development, allow users to record symptoms and appointments, and offer curated health
content. They are informative but passive: they display content without analysing the
user's own health data.

\textbf{Doctor appointment booking platforms} allow patients to search for doctors by
speciality and location, view their availability, and reserve a time slot. They excel at
scheduling but offer no health monitoring, no medication safety information, and no
management of medical records.

\textbf{Hospital information systems} centralise patient records within a single institution
and give practitioners access to a shared medical history. They are powerful within their
walls but are not accessible to patients outside those walls, and they serve institutions
rather than individuals.



\subsection{Critique of Existing Solutions}
The review of existing solutions reveals a consistent pattern: each platform addresses one
or two dimensions of the problem well, but none brings all the necessary capabilities
together in a single experience tailored to the Tunisian context.

\textbf{BabyCenter} is one of the most widely used pregnancy applications among Tunisian
women due to its availability in French and its rich content. It offers week-by-week
fetal development tracking, symptom logging, kick counting, and weight monitoring.
However, it remains an informational and community tool rather than a clinical one.
It has no connection to local doctors, no medication safety guidance, and no automated
health risk evaluation based on recorded measurements.

\textbf{Grossesse+}, equally popular among French-speaking Tunisian users, provides
detailed week-by-week pregnancy guides enriched with interactive 3D models, weight
tracking, contraction and kick counters, and appointment reminders. Like BabyCenter,
it functions purely as a reference and personal diary application with no integration
into the healthcare system and no clinical safety analysis of any kind.

\textbf{Keeplyna} is a Tunisian health platform that allows users to manage a digital
medical record including prescriptions, blood pressure readings, allergy information,
and chronic conditions, with an official list of doctors registered through the Tunisian
Medical Council. However, it is a general health management tool with no features
specific to pregnancy monitoring, no medication safety analysis for pregnant women,
and no doctor appointment scheduling functionality.

\textbf{Docic} is a Tunisian digital health platform founded in 2022 that enables patients
and healthcare professionals to store, organise, and share medical records. It supports
document sharing between patients and doctors, which is a meaningful step toward
digitising patient records in Tunisia. However, it has no pregnancy-specific features,
no health measurement recording, no risk evaluation engine, and no medication
safety guidance.

Across all four platforms, four critical capabilities are universally absent: automatic
evaluation of health risk based on recorded measurements, medication safety information
adapted to pregnancy and the Tunisian drug market, a doctor appointment scheduling
system integrated with health monitoring, and an interface specifically designed for
pregnant women.
These four gaps define exactly what Sahhty was built to fill.

% ── 1.6 Proposed Solution ─────────────────────────────────────
\section{Proposed Solution}

Sahhty is a mobile health platform designed to give pregnant women and all patients a
single, coherent tool to manage their health throughout their care journey. The platform
brings together six interconnected capabilities:

\textbf{Automatic health risk evaluation.} Each time the patient records a new health
reading, the platform automatically evaluates the result against the patient's full history
and produces an assessment of the overall risk level. If the result signals a concern, the
patient receives an immediate alert on her screen and a recommendation to consult her
doctor. High-concern results also generate a notification that reaches the patient even
when the application is closed.

\textbf{Complete health monitoring.} Patients record six types of health readings: weight,
blood pressure, blood sugar level, heart rate, body temperature, and body mass index. The
full history of readings is available at any time, with the ability to filter by type, sort by
date, and navigate through past entries.

\textbf{Pregnancy follow-up.} The application tracks the pregnancy week by week from
the day it is discovered until the expected delivery date. A visual progress indicator and a
fruit size comparison help the patient understand each stage of development intuitively.
For women, the platform also tracks the menstrual cycle. Male patients see a completely
different interface that shows only the information relevant to their situation.

\textbf{Doctor search and appointment booking.} Patients search for doctors by speciality,
city, or distance on an interactive map. Each doctor profile shows their schedule, years
of experience, and consultation fees. Appointments are booked directly from the
application, and both the patient and the doctor receive an immediate notification when a
booking is made or confirmed.

\textbf{Personal medical file management.} Patients upload and organise their health
documents by category. They decide which doctors are allowed to view their files. A doctor
who has been granted access can view all documents, alongside the patient's recent health
readings and current treatments. Access can be revoked at any time by the patient.

\textbf{Medication safety and drug interaction checking.} The platform integrates the full
list of medications available in Tunisia, enriched with pregnancy safety information
sourced from reference databases used by healthcare professionals. For each medication,
the patient can see whether it is considered safe, to be used with caution, or to be avoided
at each stage of pregnancy. The platform also warns the patient if a medication contains a
substance she is known to be allergic to, or if two medications in her treatment plan are
known to interact with each other.

\textbf{Smartwatch data collection.} A companion application installed on a paired
smartwatch automatically measures the patient's heart rate every thirty minutes and sends
the reading to the phone, which then forwards it to the platform. This allows health data
to be collected passively without requiring the patient to do anything.

The platform also sends instant notifications for critical events, supports two languages
(French and English), and provides a completely distinct visual experience for male patients
compared to female patients, adapting the colors and content to each user's profile.

% ── 1.7 Methodology: V-Model ─────────────────────────────────
\section{Methodology: V-Model}

\subsection*{Presentation of the V-Model}

The development of Sahhty was guided by the V-Model, a structured approach to software
development in which every stage of construction is directly paired with a corresponding
stage of verification. The model takes its name from the shape it forms when represented
visually: the left side descends through a series of definition stages, and the right side
rises through a corresponding series of testing stages, meeting at the point where the
system is built.As shown in (Figure 1.1) .

\begin{figure}[H]
    \centering
    \includegraphics[width=0.75\textwidth]{images/vmodel.png}
    \caption{The V-Model development approach applied to the Sahhty project}
    \label{fig:vmodel}
\end{figure}

The V-Model was chosen for this project for two main reasons. First, it enforces a
discipline of thinking about how each piece of the system will be verified \textit{before}
it is built, which reduces the risk of discovering fundamental problems late in the
development process. Second, it provides clear milestones and deliverables at each stage,
making it easier for a two-person team to stay aligned and to know when a phase is
genuinely complete before moving to the next one.

The following (Table 1.1) describes how each stage of the V-Model was applied to the Sahhty
project.

\begin{table}[H]
\centering
\caption{Application of the V-Model stages to the Sahhty project}
\renewcommand{\arraystretch}{1.45}
\begin{tabular}{|p{4.5cm}|p{8.5cm}|}
\hline
\textbf{Stage} & \textbf{What was done} \\
\hline
Needs analysis &
The team identified who would use the platform, what problems they face, and what the
system must do to address those problems. This stage produced the list of features and
the acceptance criteria that defined the final target. \\
\hline
Overall system design &
The main building blocks of the platform were defined: the mobile application, the
server managing the data, and the database storing all information. The way these three
blocks communicate was agreed upon before any construction began. \\
\hline
Detailed design &
Each feature was broken down into its individual parts. The screens, the data structures,
and the interactions between the phone and the server were described in detail through
diagrams. \\
\hline
Construction &
The mobile application and the server were built in parallel by the two team members,
following the designs established in the previous stages. Features were integrated
progressively and verified at each step. \\
\hline
Component verification &
Each individual feature was tested in isolation to confirm that it behaves correctly
under normal conditions and handles errors gracefully. \\
\hline
Integration verification &
The mobile application and the server were connected and tested together using real
accounts and real data to confirm that information flows correctly between them. \\
\hline
Overall system validation &
The complete platform was tested from the user's point of view, following real usage
scenarios for patients and doctors, to confirm that the system meets the needs
identified at the start. \\
\hline
Acceptance &
The platform was demonstrated on a real device to confirm that it meets the goals
defined during the internship, in terms of functionality, reliability, and ease of use. \\
\hline
\end{tabular}
\end{table}

\subsection{Gantt Chart}

The project was carried out over a period of approximately four months. The planning was
organised into successive phases corresponding to the stages of the V-Model, with the
two team members working in parallel on their respective responsibilities.As shown in (Figure 1.2) and (Table 1.2) .

\begin{figure}[H]
    \centering
    \includegraphics[width=0.98\textwidth]{images/gantt_chart.png}
    \caption{Gantt chart (Sahhty project planning and timeline)}
    \label{fig:gantt}
\end{figure}

\begin{table}[H]
\centering
\caption{Project phase summary}
\renewcommand{\arraystretch}{1.4}
\begin{tabular}{|p{5cm}|p{2.5cm}|p{5.5cm}|}
\hline
\textbf{Phase} & \textbf{Duration} & \textbf{Main outcome} \\
\hline
Needs analysis and design      & Weeks 1--3  & Features defined, diagrams produced \\
\hline
Authentication and accounts    & Weeks 4--5  & Login, registration, and identity verification working \\
\hline
Health monitoring              & Weeks 6--8  & Health readings recorded, risk evaluated automatically \\
\hline
Doctors and appointments       & Weeks 9--10 & Doctor search, map, and booking live \\
\hline
Medical files and medications  & Weeks 11--12 & Document storage, drug safety checker complete \\
\hline
Smartwatch and final features  & Weeks 13--14 & Watch data collected, male interface, two languages \\
\hline
Testing and validation         & Weeks 15--16 & Full platform verified on real device \\
\hline
Report writing                 & Weeks 14--17 & Final report submitted \\
\hline
\end{tabular}
\label{tab:timeline}
\end{table}


% ── 1.8 Conclusion ────────────────────────────────────────────
\section{Conclusion}

This chapter has laid the foundations of the Sahhty project. Starting from the observation
that pregnant women and patients in Tunisia lack a single, integrated tool to manage their
health, a clear problem statement was established around five recurring difficulties:
fragmented medical records, the absence of automatic risk evaluation, limited doctor-patient
coordination, the lack of reliable medication safety information during pregnancy, and the
absence of passive data collection through wearable devices.

A review of four representative applications confirmed that no existing solution addresses
all of these difficulties together. Each platform covers a narrow part of the problem well,
but leaves the rest unaddressed. This gap directly motivated the design of Sahhty.

The proposed platform was introduced as a unified response to all five challenges, bringing
together health monitoring, pregnancy follow-up, doctor connectivity, medical file
management, medication safety, and smartwatch integration into one coherent experience.

Finally, the V-Model development approach was presented as the guiding framework for the
project, providing a structured path from the definition of needs to the final validation of
the complete system.




% ── References ───────────────────────────────────────────────
\begin{thebibliography}{9}

\bibitem{who2016}
World Health Organization (2016).
\textit{mHealth: Use of Mobile Wireless Technologies for Public Health}.
EB/139/8.

\bibitem{who2018}
World Health Organization (2018).
\textit{mHealth: Use of Appropriate Digital Technologies for Public Health}.
A71/20.

\bibitem{milne2020}
Milne-Ives M, Lam C, De Cock C, Van Velthoven MH, Meinert E.
Mobile Apps for Health Behavior Change in Physical Activity, Diet, Drug and Alcohol
Use, and Mental Health: Systematic Review.
\textit{JMIR Mhealth Uhealth} 2020;8(3):e17046.

\end{thebibliography}